\documentclass[CJK,13pt]{beamer}
\usepackage{CJKutf8}
\usepackage{beamerthemesplit}
\usetheme{Malmoe}
\useoutertheme[footline=authortitle]{miniframes}
\usepackage{amsmath}
\usepackage{amssymb}
\usepackage{graphicx}
\usepackage{eufrak}
\usepackage{color}
\usepackage{slashed}
\usepackage{simplewick}
\usepackage{tikz}
\usepackage{tcolorbox}
\usepackage{ulem}
\graphicspath{{../figures/}}
%%figures
\def\lfig#1#2{\includegraphics[width=#1 in]{#2}}
\def\tfig#1#2{\includegraphics[height=#1 in]{#2}}
\def\addfig#1#2{\begin{center}\includegraphics[width=#1 in]{#2}\end{center}}
\def\question{\includegraphics[width=0.3in]{why.jpg}\,}
\def\answer{\includegraphics[width=0.3in]{baozou_haha.png}\,}
\def\wulian{\includegraphics[width=0.18in]{emoji_wulian.jpg}}
\def\bigwulian{\includegraphics[width=0.35in]{emoji_wulian.jpg}}
\def\bye{\includegraphics[width=0.18in]{emoji_bye.jpg}}
\def\bigbye{\includegraphics[width=0.35in]{emoji_bye.jpg}}
\def\huaixiao{\includegraphics[width=0.18in]{emoji_huaixiao.jpg}}
\def\bighuaixiao{\includegraphics[width=0.35in]{emoji_huaixiao.jpg}}
\def\jianxiao{\includegraphics[width=0.18in]{emoji_jianxiao.jpg}}
\def\bigjianxiao{\includegraphics[width=0.35in]{emoji_jianxiao.jpg}}
\def\haoqi{\includegraphics[width=0.18in]{emoji_haoqi.jpg}}
%% colors
\def\blacktext#1{{\color{black}#1}}
\def\bluetext#1{{\color{blue}#1}}
\def\redtext#1{{\color{red}#1}}
\def\darkbluetext#1{{\color[rgb]{0,0.2,0.6}#1}}
\def\skybluetext#1{{\color[rgb]{0.2,0.7,1.}#1}}
\def\cyantext#1{{\color[rgb]{0.,0.5,0.5}#1}}
\def\greentext#1{{\color[rgb]{0,0.7,0.1}#1}}
\def\darkgray{\color[rgb]{0.2,0.2,0.2}}
\def\lightgray{\color[rgb]{0.6,0.6,0.6}}
\def\gray{\color[rgb]{0.4,0.4,0.4}}
\def\blue{\color{blue}}
\def\red{\color{red}}
\def\orange{\color[rgb]{1.,0.8,0.}}
\def\green{\color{green}}
\def\darkgreen{\color[rgb]{0,0.4,0.1}}
\def\darkblue{\color[rgb]{0,0.2,0.6}}
\def\skyblue{\color[rgb]{0.2,0.7,1.}}
%%control
\def\diag{\mathrm{diag}\,}
\def\heaviside{\,\mathrm{h}\,}
\def\bral{\left(\begin{array}{l}}
\def\brar{\end{array}\right)}
\def\brall{\left(\begin{array}{ll}}
\def\brarr{\end{array}\right)}
\def\bralll{\left(\begin{array}{lll}}
\def\brarrr{\end{array}\right)}
\def\branchl{\left\{\begin{array}{l}}
\def\branchr{\end{array}\right.}
\def\branchll{\left\{\begin{array}{ll}}
\def\branchrr{\end{array}\right.}
\def\branchlll{\left\{\begin{array}{lll}}
\def\branchrrr{\end{array}\right.}
\def\sfgamma#1{\,\Gamma\left( #1 \right)\,}
\def\be{\begin{equation}}
\def\ee{\nonumber\end{equation}}
\def\bea{\begin{eqnarray}}
\def\eea{\nonumber\end{eqnarray}}
\def\bch{\begin{CJK}{UTF8}{gbsn}}
\def\ech{\end{CJK}}
\def\bitem{\begin{itemize}}
\def\eitem{\end{itemize}}
\def\bcenter{\begin{center}}
\def\ecenter{\end{center}}
\def\bex{\begin{minipage}{0.2\textwidth}\includegraphics[width=0.6in]{jugelizi.png}\end{minipage}\begin{minipage}{0.76\textwidth}}
\def\eex{\end{minipage}}
\def\chtitle#1{\frametitle{\bch#1\ech}}
\def\bmat#1{\left(\begin{array}{#1}}
\def\emat{\end{array}\right)}
\def\bcase#1{\left\{\begin{array}{#1}}
\def\ecase{\end{array}\right.}
\def\bmini#1{\begin{minipage}{#1\textwidth}}
\def\emini{\end{minipage}}
\def\tbox#1{\begin{tcolorbox}#1\end{tcolorbox}}
\def\pfrac#1#2#3{\left(\frac{\partial #1}{\partial #2}\right)_{#3}}
\def\res#1#2{\,\mathrm{res}\,#1\left(#2\right)\,}
\def\newt#1#2{\left(\begin{array}{c}#1\\ #2\end{array}\right)}
\def\reof#1{\,\mathrm{Re}{\left(#1\right)\,}}
\def\imof#1{\,\mathrm{Im}{\left(#1\right)\,}}
\def\Arg#1{\,\mathrm{Arg}\,#1\,}
%%symbols
\def\intfull{\int_{-\infty}^\infty}
\def\inthalf{\int_0^\infty}
\def\bropt{\,(\ \ \ )}
\def\sech{\mathrm{sech}\,}
\def\csch{\mathrm{csch}\,}
\def\asinh{\mathrm{asinh}\,}
\def\acosh{\mathrm{acosh}\,}
\def\sone{$\star$}
\def\stwo{$\star\star$}
\def\sthree{$\star\star\star$}
\def\sfour{$\star\star\star\star$}
\def\sfive{$\star\star\star\star\star$}
\def\rint{{\int_\leftrightarrow}}
\def\roint{{\oint_\leftrightarrow}}
\def\stdHf{{\textit{\r H}_f}}
\def\deltaH{{\Delta \textit{\r H}}}
\def\ii{{\dot{\imath}}}
\def\skipline{{\vskip0.1in}}
\def\skiplines{{\vskip0.2in}}
\def\lagr{{\mathcal{L}}}
\def\hamil{{\mathcal{H}}}
\def\vecT{{\mathbf{T}}}
\def\vecN{{\mathbf{N}}}
\def\vecB{{\mathbf{B}}}
\def\vecv{{\mathbf{v}}}
\def\vecr{{\mathbf{r}}}
\def\vecf{{\mathbf{f}}}
\def\vecg{{\mathbf{g}}}
\def\vecupsilon{{\mathbf{\upsilon}}}
\def\vecu{{\mathbf{u}}}
\def\vecj{{\mathbf{j}}}
\def\vecx{{\mathbf{x}}}
\def\vecy{{\mathbf{y}}}
\def\vecz{{\mathbf{z}}}
\def\veck{{\mathbf{k}}}
\def\vecp{{\mathbf{p}}}
\def\vecn{{\mathbf{n}}}
\def\vecA{{\mathbf{A}}}
\def\vecP{{\mathbf{P}}}
\def\vecsigma{{\mathbf{\sigma}}}
\def\hatJn{{\hat{J_\vecn}}}
\def\hatJx{{\hat{J_x}}}
\def\hatJy{{\hat{J_y}}}
\def\hatJz{{\hat{J_z}}}
\def\hatj#1{\hat{J_{#1}}}
\def\hatphi{{\hat{\phi}}}
\def\hatq{{\hat{q}}}
\def\hatpi{{\hat{\pi}}}
\def\vel{\upsilon}
\def\Dint{{\mathcal{D}}}
\def\adag{{\hat{a}^\dagger}}
\def\bdag{{\hat{b}^\dagger}}
\def\cdag{{\hat{c}^\dagger}}
\def\ddag{{\hat{d}^\dagger}}
\def\hata{{\hat{a}}}
\def\hatb{{\hat{b}}}
\def\hatc{{\hat{c}}}
\def\hatd{{\hat{d}}}
\def\hatD{{\,\hat{D}}}
\def\hatN{{\hat{N}}}
\def\hatH{{\hat{H}}}
\def\hatp{{\hat{p}}}
\def\Fup{{F^{\mu\nu}}}
\def\Fdown{{F_{\mu\nu}}}
\def\newl{\nonumber \\}
\def\vece{\mathrm{e}}
\def\calM{{\mathcal{M}}}
\def\calT{{\mathcal{T}}}
\def\calR{{\mathcal{R}}}
\def\barpsi{\bar{\psi}}
\def\baru{\bar{u}}
\def\barv{\bar{\upsilon}}
\def\qeq{\stackrel{?}{=}}
\def\ftf{\stackrel{\mathcal{FT}}{\Longrightarrow}}
\def\ftb{\stackrel{\mathcal{FT}}{\Longleftarrow}}
\def\ftfb{\stackrel{\mathcal{FT}}{\Longleftrightarrow}}
\def\ltf{\stackrel{\mathcal{LT}}{\Longrightarrow}}
\def\ltb{\stackrel{\mathcal{LT}}{\Longleftarrow}}
\def\ltfb{\stackrel{\mathcal{LT}}{\Longleftrightarrow}}
\def\torder#1{\mathcal{T}\left(#1\right)}
\def\rorder#1{\mathcal{R}\left(#1\right)}
\def\contr#1#2{\contraction{}{#1}{}{#2}#1#2}
\def\trof#1{\mathrm{Tr}\left(#1\right)}
\def\trace{\mathrm{Tr}}
\def\comm#1{\ \ \ \left(\mathrm{used}\ #1\right)}
\def\tcomm#1{\ \ \ (\text{#1})}
\def\slp{\slashed{p}}
\def\slk{\slashed{k}}
\def\calp{{\mathfrak{p}}}
\def\veccalp{\mathbf{\mathfrak{p}}}
\def\Tthree{T_{\tiny \textcircled{3}}}
\def\pthree{p_{\tiny \textcircled{3}}}
\def\dbar{{\,\mathchar'26\mkern-12mu d}}
\def\erf{\mathrm{erf}}
\def\const{\mathrm{const.}}
\def\pheat{\pfrac p{\ln T}V}
\def\vheat{\pfrac V{\ln T}p}
%%units
\def\fdeg{{^\circ \mathrm{F}}}
\def\cdeg{^\circ \mathrm{C}}
\def\atm{\,\mathrm{atm}}
\def\angstrom{\,\text{\AA}}
\def\SIL{\,\mathrm{L}}
\def\SIkm{\,\mathrm{km}}
\def\SIyr{\,\mathrm{yr}}
\def\SIGyr{\,\mathrm{Gyr}}
\def\SIV{\,\mathrm{V}}
\def\SImV{\,\mathrm{mV}}
\def\SIeV{\,\mathrm{eV}}
\def\SIkeV{\,\mathrm{keV}}
\def\SIMeV{\,\mathrm{MeV}}
\def\SIGeV{\,\mathrm{GeV}}
\def\SIcal{\,\mathrm{cal}}
\def\SIkcal{\,\mathrm{kcal}}
\def\SImol{\,\mathrm{mol}}
\def\SIN{\,\mathrm{N}}
\def\SIHz{\,\mathrm{Hz}}
\def\SIm{\,\mathrm{m}}
\def\SIcm{\,\mathrm{cm}}
\def\SIfm{\,\mathrm{fm}}
\def\SImm{\,\mathrm{mm}}
\def\SInm{\,\mathrm{nm}}
\def\SImum{\,\mathrm{\mu m}}
\def\SIJ{\,\mathrm{J}}
\def\SIW{\,\mathrm{W}}
\def\SIkJ{\,\mathrm{kJ}}
\def\SIs{\,\mathrm{s}}
\def\SIkg{\,\mathrm{kg}}
\def\SIg{\,\mathrm{g}}
\def\SIK{\,\mathrm{K}}
\def\SImmHg{\,\mathrm{mmHg}}
\def\SIPa{\,\mathrm{Pa}}
%page
\def\secpage#1#2{\begin{frame}\bcenter{\bf \Huge #1} \ecenter \skipline \tbox{\bcenter #2 \ecenter}\end{frame}}
\def\append#1#2{\secpage{附录#1}{#2}}
\def\thinka#1{\begin{frame}\frametitle{思考题}\bcenter\lfig{0.4}{think0.jpg}  \skipline #1 \ecenter\end{frame}}
\def\thinkb#1{\begin{frame}\frametitle{思考题}\bcenter\lfig{0.5}{think1.jpg}  \skipline #1 \ecenter\end{frame}}
\def\thinkc#1{\begin{frame}\frametitle{思考题}\bcenter\lfig{0.5}{think2.jpg}  \skipline #1 \ecenter\end{frame}}
\def\thinkd#1{\begin{frame}\frametitle{思考题}\bcenter\lfig{0.5}{think3.jpg}  \skipline #1 \ecenter\end{frame}}
\def\thinke#1{\begin{frame}\frametitle{思考题}\bcenter\lfig{0.5}{think4.jpg}  \skipline #1 \ecenter\end{frame}}
\def\thinkf#1{\begin{frame}\frametitle{思考题}\bcenter\lfig{0.5}{think5.jpg}  \skipline #1 \ecenter\end{frame}}
\def\schw{ds^2 = \left(1-\frac{2GM}{r}\right)dt^2 - \left(1-\frac{2GM}{r}\right)^{-1} dr^2 - r^2\left(d\theta^2 + \sin^2\theta d\phi^2\right)}

\input{titlepage.tex}
  \date{}
  \begin{document}
  \bch
  \tpage{22}{Gravitational Radiation (I)}

  \begin{frame}
    之前已经给出了引力波的推迟势解，在这一讲我们要讲如何实际应用它。

    \skiplines

    为了让讨论清晰简洁，我们有必要引入一个非常好用的操作——

    \addfig{1.2}{kaishicaozuo.jpg}
    
  \end{frame}
  
  \begin{frame}
    \frametitle{bar操作}
    对任何一个线性化的，用 $\eta_{\mu\nu}$ 进行指标升降的小量 $f_{\mu\nu}$，我们都可以定义
    $$\bar{f}_{\mu\nu} \equiv f_{\mu\nu}-\frac{1}{2}\eta_{\mu\nu} f^{\alpha}_{\ \alpha}$$
    那么线性化的爱因斯坦方程就可以写成
      $$ \bar{R}_{\mu\nu} = 8\pi GT_{\mu\nu}  $$

    \skipline

    更有意思的是，两次加bar操作相当于没有操作，也就是爱因斯坦方程两边加bar可以得到
    $$R_{\mu\nu} = 8\pi G\bar{T}_{\mu\nu}  $$
  \end{frame}


  \begin{frame}
    \frametitle{bar操作(续）}
      在谐和坐标系里，$R_{\mu\nu}=-\frac{1}{2}\square h_{\mu\nu}$。爱因斯坦方程成为
      $\square h_{\mu\nu} = -16\pi G \bar{T}_{\mu\nu}$
      或者等价的
      $\square \bar{h}_{\mu\nu} = -16\pi G T_{\mu\nu}$。

      
      因此推迟势解可以写成我们之前给出过的：
      {\blue $$ h_{\mu\nu}(t, \vecx) = -4G \int d^3\vecx' \frac{\bar{T}_{\mu\nu}\left(\vecx', t - \lvert \vecx -\vecx'\rvert\right)}{\lvert \vecx - \vecx'\rvert} $$}
      或者其等价形式
      {\blue $$ \bar{h}^{\mu\nu}(t, \vecx) = -4G \int d^3\vecx' \frac{T^{\mu\nu}\left(\vecx', t - \lvert \vecx -\vecx'\rvert\right)}{\lvert \vecx - \vecx'\rvert} $$}
      
  \end{frame}


  \secpage{引力波的能量密度}{$$ \rho_{\rm gw} = \frac{1}{32\pi G}\left(\dot h_+^2 + \dot h_\times^2 + |\nabla h_+|^2 + |\nabla h_\times|^2\right)$$}


  \begin{frame}
    先给个可能会令你失望的结论：{\blue 引力没有真正意义上的局域的能量动量张量}，也就是你问在这个时空点的引力能量密度是多少，一般是没有明确答案的。
    (可以参考 Dirac 或者 Weinberg 书上的讨论)。

    \skipline

    {\scriptsize 这决不仅仅是因为引力的表述形式里具有非物理的自由度（电磁场的四维矢势里也有，但电磁场明确地具有局域的能量动量张量）。首先，非物理的自由度和张量的定义（“坐标变换”）纠缠到了一起，把问题变得更复杂。其次，引力的4个非内禀自由度只是被动地响应物质的分布——这种没有决定自己命运的能力的“傀儡”是否应该具有“能量动量”，本身也是一件值得商榷的事情。}
    
    \skipline

    不过，对{\blue 单方向传播的引力波，我们可以操作掉所有非内禀的自由度，计算它的“物理”能量密度}。

    \skipline
    
    {\scriptsize 对于多方向传播的引力波交杂在一起的情况，我们并没有统一的“操作坐标系”的方式来消除所有非内禀自由度，把各种不同坐标系的能量密度加起来的操作就有点诡异了（不过我的个人观点是这问题不大）。}
    
  \end{frame}
  
  \begin{frame}
    \frametitle{逆变度规的变化量}
    如果我们约定用 $\eta_{\mu\nu}$ 对 $h_{\mu\nu}$ 进行指标升降，那么容易验证
    $\eta^{\mu\nu}-h^{\mu\nu}$ 和 $\eta_{\mu\nu}+h_{\mu\nu}$ 在一阶近似意义下互为逆矩阵。

    \skiplines
    
    也就是说 $-h^{\mu\nu}$ 是对 Minkowski 逆变形式度规的偏离。
  \end{frame}

  \begin{frame}
    \frametitle{近似作用量}
    在讨论变分法推导爱因斯坦方程时我们证明了：
    $$ \delta S = -\frac{1}{16\pi G}\int \sqrt{-g} d^4x \delta g^{\mu\nu} \left(R_{\mu\nu}-\frac{1}{2}R g_{\mu\nu}\right) $$
    如果从 Minkowski 度规出发，微扰到 $\eta_{\mu\nu}+h_{\mu\nu}$，那么由于Minkowski度规给出的 $R_{\mu\nu}=0$，上面的积分也是零。

    \skipline

    \bmini{0.55}
    这并不奇怪，说的无非是当没有物质作用量时，Minkowski度规让作用量取到稳定值（一阶变化量消失）。
    
    所以我们要算的是个二阶小量——
    \emini
    \bmini{0.4}
    \addfig{1.}{youkaishicaozuo.jpg}
    \emini
  \end{frame}


  \begin{frame}
    \frametitle{近似作用量（续）}
    把度规 $\eta_{\mu\nu}+\epsilon h_{\mu\nu}$ 微扰到  $\eta_{\mu\nu}+(\epsilon+d\epsilon)h_{\mu\nu}$ 时，
    $$ \delta g^{\mu\nu} \approx -h^{\mu\nu}d\epsilon,\ \ R_{\mu\nu}-\frac{1}{2}Rg_{\mu\nu} \approx -\frac{1}{2}\epsilon \square \bar{h}_{\mu\nu} $$
    这里利用前一讲推导的谐和坐标系里的 $R_{\mu\nu}\approx -\frac{1}{2}\square h_{\mu\nu}$。
    
    所以
    \bea
    dS &=& S(\epsilon+d\epsilon)-S(\epsilon) \newl
    &=& -\frac{1}{16\pi G}\int \sqrt{-g} d^4x\, \left(-d\epsilon \, h^{\mu\nu}\right)\left(-\frac{1}{2}\epsilon \square \bar{h}_{\mu\nu}\right) \newl
    &\approx & -\frac{\epsilon d\epsilon}{32\pi G}\int \sqrt{-g}  d^4x \,  h^{\mu\nu}\square \bar{h}_{\mu\nu}
    \eea
  \end{frame}

  

  \begin{frame}
    \frametitle{近似作用量（续）}
    然后，很简单，近似把 $\sqrt{-g}$ 当成 $1$，并用普通微商代替协变微商 (因为已经是二阶小量了，再修正下意义不大)。把 $\epsilon$ 从 $0$ 积分到 $1$，就会得到
    \bea
    S &=& -\frac{1}{64\pi G} \int d^4x\, h^{\mu\nu}\partial^\alpha\partial_\alpha \bar{h}_{\mu\nu} \newl
    &=& \frac{1}{64\pi G} \int d^4x\, \partial^\alpha h^{\mu\nu} \partial_\alpha \bar{h}_{\mu\nu} + \text{ignored surface terms}
    \eea
    于是我们得到了作用量密度
    \be
    \mathcal{L} = \frac{1}{64\pi G}\left(\frac{\partial h^{\mu\nu}}{\partial t}\frac{\partial \bar{h}_{\mu\nu}}{\partial t} - \nabla h^{\mu\nu}\cdot \nabla \bar{h}_{\mu\nu}\right) 
    \ee
  \end{frame}

  \begin{frame}
    \frametitle{哈密顿量密度}
    对应 $h_{\mu\nu}$ 的广义动量
    $$ p^{\mu\nu} = \frac{\partial \mathcal{L}}{\partial \left(\frac{\partial h_{\mu\nu}}{\partial t}\right)} = \frac{1}{32\pi G}\frac{\partial \bar{h}^{\mu\nu}}{\partial t}$$
    对应的“哈密顿量密度（能量密度）”
    {\blue $$ \mathcal{H} = \frac{\partial h_{\mu\nu}}{\partial t} p^{\mu\nu} - \mathcal{L}  = \frac{1}{64\pi G}\left(\frac{\partial \bar{h}^{\mu\nu}}{\partial t}\frac{\partial h_{\mu\nu}}{\partial t} + \nabla \bar{h}^{\mu\nu}\cdot \nabla h_{\mu\nu}\right) $$}
    如前所述，这个“能量密度”并不是通常意义上的局域能量动量张量的00分量，对其物理意义进行解释时须谨慎！
  \end{frame}
  

  \begin{frame}
    考虑一个平面波，取其传播方向为 $z$ 轴，就有
    $$ h_{11} = h_+, h_{22} = -h_+, h_{12} = h_{21}= h_\times $$
    其余分量均为零。

    \skipline

    容易算出
    $$ \bar{h}^{11} = h_+, \bar{h}^{22} = -h_+, \bar{h}^{12} = \bar{h}^{21}= h_\times $$
    则能量密度{\blue
      $$ \rho_{\rm gw} = \frac{1}{32\pi G}\left(\dot h_+^2 + \dot h_\times^2 + |\nabla h_+|^2 + |\nabla h_\times|^2\right)$$}
  \end{frame}

  \secpage{引力辐射}{只是个傅立叶变换而已！}

  \begin{frame}
    考虑以频率$\omega$变化，分布在有限区域内的引力波源
    $$T_{\mu\nu}(t, \vecx) = M_{\mu\nu}(\vecx)e^{-i\omega t} +c.c.$$
    {\scriptsize $c.c.$ 表示复共轭}
      
    在离波源很远的 $\vecx = r\vecn $ ($\vecn$ 是代表方向的单位矢, $r$ 代表距离) 处观测，由之前推出的辐射公式有
    $$ \bar{h}_{\mu\nu}(t, \vecx) = -4G \int \frac{d^3\vecx'}{\lvert \vecx - \vecx' \rvert}M_{\mu\nu}(\vecx')e^{-i\omega \left(t - \lvert\vecx-\vecx'\rvert\right)} +c.c.$$


    \bmini{0.5}
    \lfig{1.2}{gwcalc.png}
    \emini
    \bmini{0.45}

    \bea
    |\vecx-\vecx'|&=& XX'\newl
    &\approx & XA \newl
    &=& OX - OA \newl
    &=& r - \vecx'\cdot \vecn
    \eea
    \emini
  \end{frame}

  \begin{frame}
    把振幅取零阶近似，相位需要更精确些，取一阶近似：
     \bea
    \bar{h}_{\mu\nu}(t, \vecx) &=& -4G \int \frac{d^3\vecx'}{r}M_{\mu\nu}(\vecx')e^{-i\omega \left(t - r+ \vecx'\cdot \vecn\right)} +c.c.\newl
    &\equiv & -\frac{4G e^{-i\omega(t-r)}}{r}\mathcal{M}_{\mu\nu}(\omega\vecn) +c.c.
    \eea
    这里的
    $$\mathcal{M}_{\mu\nu}(\veck)\equiv \int d^3\vecx M_{\mu\nu}(\vecx) e^{-i\veck\cdot\vecx}$$
    是 $M_{\mu\nu}(\vecx)$ （$T_{\mu\nu}$ 的振幅）的三维傅立叶变换。
  \end{frame}


  \begin{frame}
    通过把结果两边取 bar 以及升指标，我们得到
    {\blue $$ h_{\mu\nu}(t, r\vecn) = - \frac{4G e^{-i\omega(t-r)}}{r} \bar{\mathcal{M}}_{\mu\nu}(\omega\vecn) + c.c.$$}
    以及等价地
    {\blue $$ \bar{h}^{\mu\nu}(t, r\vecn) = - \frac{4G e^{-i\omega(t-r)}}{r} \mathcal{M}^{\mu\nu}(\omega\vecn) + c.c.$$}
  \end{frame}
  
  \begin{frame}
    在 $\vecn$ 方向附近的 $d\Omega$ 的立体角内，距离在 $r$ 和 $r+\delta t$ 之间的空间范围内的引力波总能量，显然是源在 $\delta$ 时间间隔内发出的总能量。于是有
    \be
    \frac{dP}{d\Omega} = \rho_{\rm gw}(t, r\vecn) r^2
    \ee
    我们在远大于 $\frac{1}{\omega}$ 的时间内对它求平均(用 $\langle\ldots\rangle$表示)，只有相位抵消的项才有贡献，所以我们可以把 $\partial/\partial t$ 和 $\nabla$  都替换为 $\omega$ (实际是 $i\omega$ 或者 $-i\omega$，但有非零贡献的项总是一个 $i\omega$ 和 一个 $-i\omega$ 相乘)：
    \bea
    \left\langle \frac{dP}{d\Omega} \right\rangle &=&  \frac{r^2}{64\pi G}\left\langle\frac{\partial \bar{h}^{\mu\nu}}{\partial t}\frac{\partial h_{\mu\nu}}{\partial t} + \nabla \bar{h}^{\mu\nu}\cdot \nabla h_{\mu\nu}\right\rangle \newl
    &=& \frac{r^2\omega^2}{32\pi G} \left\langle \bar{h}^{\mu\nu}h_{\mu\nu} \right\rangle
    \eea
    这里我们对 $h_{\mu\nu}$ 都省略了宗量 $(t, r\vecn)$。
    \end{frame}

  \begin{frame}
    代入前面求解得的 $\bar{h}^{\mu\nu}$ 和 $h_{\mu\nu}$
    {\small
    \bea
    \left\langle \frac{dP}{d\Omega} \right\rangle &=&   \frac{G \omega^2}{2\pi } \left\langle \left(e^{-i\omega(t-r)} \mathcal{M}^{\mu\nu}(\omega\vecn)+c.c.\right)  \left(e^{-i\omega(t-r)} \bar{\mathcal{M}}_{\mu\nu}(\omega\vecn)+c.c.\right) \right\rangle \newl
    &=&   \frac{G \omega^2}{2\pi }\left(\mathcal{M}^{\mu\nu}(\omega\vecn)\bar{\mathcal{M}}_{\mu\nu}(\omega\vecn)^* + c.c.\right) \newl
    &=&   \frac{G \omega^2}{\pi }\left(\mathcal{M}^{\mu\nu}(\omega\vecn)\mathcal{M}_{\mu\nu}^*(\omega\vecn) -\frac{1}{2}\lvert \mathcal{M}^\alpha_{\ \alpha}(\omega\vecn)\rvert^2\right)    
    \eea    
    }
    通常我们会省略取平均的符号（因为这是显然必须做的）——
  \end{frame}

  \begin{frame}
    \frametitle{单频引力辐射功率公式}
    \tbox{
      对单频率源，在方向 $\vecn$，单位立体角内的引力辐射功率为
      $$\frac{dP}{d\Omega} =  \frac{G \omega^2}{\pi }\left(\mathcal{M}^{\mu\nu}(\omega\vecn)\mathcal{M}_{\mu\nu}^*(\omega\vecn) -\frac{1}{2}\lvert \mathcal{M}^\alpha_{\ \alpha}(\omega\vecn)\rvert^2\right)$$
      这里的 $\mathcal{M}_{\mu\nu}$ 是 $T_{\mu\nu}$ 的振幅的三维傅立叶变换。
      }
  \end{frame}

  \secpage{连续频谱的辐射总能量}{四维傅立叶变换}
  
  \begin{frame}
    \frametitle{连续频谱}
    如果有多个单频的源，只要在足够长的时间内求平均，所有干涉项都会消失，你只要把单独每个源的辐射功率加起来就行了。

    \skiplines

    如果是在一段有限时间内发生的事件造成的连续频谱，则只能讨论辐射总能量 (因为功率不是恒定的且如何平均的意义不明确)。我们来详细讨论下——
  \end{frame}


  \begin{frame}
    考虑到所有频率的贡献，$h_{\mu\nu}$ 的解要修改成
    {\blue $$ h_{\mu\nu}(t, r\vecn) = - \int_{-\infty}^\infty d\omega \frac{2G e^{-i\omega(t-r)}}{\pi r} \bar{\mathcal{T}}_{\mu\nu}(k)$$}
    注意我用负频率 $\omega<0$ 部分的贡献替代了原先写作 $c.c.$ 的部分。这里的 $k$ 代表四维波矢 $k^\mu = (\omega, \omega\vecn)$。$\mathcal{T}_{\mu\nu}$是 $T_{\mu\nu}$ 的四维傅立叶变换：
    $$\mathcal{T}_{\mu\nu}(k) = \int d^4x \, T_{\mu\nu}(x) e^{ik_\alpha x^\alpha} $$
    \skipline

    你当然也不难写出
    {\blue $$ \bar{h}^{\mu\nu}(t, r\vecn) = - \int_{-\infty}^\infty d\omega \frac{2G e^{-i\omega(t-r)}}{\pi r} \mathcal{T}^{\mu\nu}(k)$$}
  \end{frame}
  

  \begin{frame}
    在 $\vecn$ 方向，单位立体角内的辐射总能量为
    {\scriptsize
    \bea
    \frac{dE}{d\Omega}(\vecn) &=& \frac{r^2}{64\pi G}\int_{-\infty}^\infty dt \left(\frac{\partial \bar{h}^{\mu\nu}}{\partial t}\frac{\partial h_{\mu\nu}}{\partial t} + \nabla \bar{h}^{\mu\nu}\cdot \nabla h_{\mu\nu}\right) \newl
    &=& \frac{G}{16\pi^3}\int_{-\infty}^\infty dt \int_{-\infty}^\infty d\omega \int_{-\infty}^\infty d\omega'(-2\omega\omega') e^{-i(\omega+\omega')(t-r)}\bar{\mathcal{T}}_{\mu\nu}(k)\mathcal{T}^{\mu\nu}(k') \newl
    &=& \frac{G}{8\pi^2} \int_{-\infty}^\infty d\omega \int_{-\infty}^\infty d\omega'(-2\omega\omega') \delta(\omega+\omega')\bar{\mathcal{T}}_{\mu\nu}(k)\mathcal{T}^{\mu\nu}(k') \newl
    &=& \frac{G}{4\pi^2} \int_{-\infty}^\infty d\omega \,\omega^2\bar{\mathcal{T}}_{\mu\nu}(k)\mathcal{T}^{\mu\nu}(k)^* \newl
    &=& \frac{G}{4\pi^2} \int_0^\infty d\omega \,\omega^2\left(\mathcal{T}^*_{\mu\nu}(k)\mathcal{T}^{\mu\nu}(k)-\frac{1}{2}\lvert \mathcal{T}^\alpha_{\ \alpha}(k)\rvert^2\right)    
    \eea
    }
    这个结果也可以理解为
  \end{frame}

   
  \begin{frame}
    \frametitle{连续谱引力辐射能量公式}
    \tbox{
      对连续谱源，在方向 $\vecn$，单位立体角内，单位频率间隔的引力辐射总能量为
      $$\frac{dE}{d\Omega d\omega}\left(\vecn\right) =  \frac{G\omega^2}{4\pi^2}\left(\mathcal{T}^*_{\mu\nu}(k)\mathcal{T}^{\mu\nu}(k)-\frac{1}{2}\lvert \mathcal{T}^\alpha_{\ \alpha}(k)\rvert^2\right)$$    
      这里的 $\mathcal{T}_{\mu\nu}$ 是 $T_{\mu\nu}$ 的四维傅立叶变换，四维波矢 $k^\mu=(\omega, \omega\vecn)$。
      }
  \end{frame}
  
    \ech
\end{document} 




  
